% GENERATED FILE. Edit canonical lessons, then run npm run generate:books.
% canonical-source-hash: fnv1a64:132aca507287e2e2
% canonical-lessons: ML-C260-kaccavatam, ML-C260-kaccavatakkaran, ML-C260-tulasu, ML-C260-kilo, ML-C260-sadhanam

\chapter{Business, Shopkeeper, Weighing scale, Kilo, Thing}
\label{ch:ml-business-shopkeeper-weighing-scale-kilo-thing}
\hlchaptermodality{260}

\begin{chapteropening}
\textbf{By the end of this chapter:} \emph{I can say business, a shopkeeper, a weighing scale, a kilo and a thing.}

Complete the last lesson of chapter 260: I can say business, a shopkeeper, a weighing scale, a kilo and a thing.
\end{chapteropening}

\section[\texorpdfstring{\ml{കച്ചവടം} (kaccavaṭaṁ)}{കച്ചവടം (kaccavaṭaṁ)}]{\ml{കച്ചവടം} (kaccavaṭaṁ) — business, trade}

\label{lesson:ML-C260-kaccavatam}



\begin{quote}
Before the new one: say the Malayalam for maize, then the Malayalam for chickpeas.
\end{quote}

\subsection*{You'll want to know: \ml{കച്ചവടം}}
\textbf{\ml{കച്ചവടം}} — \emph{kaccavaṭaṁ} — \textquotedblleft{}business, trade\textquotedblright{}.

\begin{cousinweb}
\textbf{\ml{കച്ചവടം}} (\emph{kaccavaṭaṁ}) is of uncertain origin; neither DEDR nor CDIAL lists it.
\end{cousinweb}

\begin{grammarlens}[title={What you've built}]
One more word for this chapter.
\end{grammarlens}

\subsection*{Your turn}
\begin{itemize}
\raggedright
  \item \emph{Say it:} \emph{kaccavaṭaṁ}
  \item \emph{Say it:} \emph{kaccavaṭaṁ}, once more
  \item \emph{Say it:} say \emph{kaccavaṭaṁ}
  \item \emph{Recall:} say the Malayalam for maize, then the Malayalam for chickpeas, then say \emph{kaccavaṭaṁ} again
\end{itemize}

\begin{tcolorbox}[breakable,colback=teal!4,colframe=teal!35!black,title={Before you move on}]
What does \ml{കച്ചവടം} mean? (\textquotedblleft{}Business, trade\textquotedblright{}.) Say it once more.
\end{tcolorbox}
\section[\texorpdfstring{\ml{കച്ചവടക്കാരൻ} (kaccavaṭakkāran)}{കച്ചവടക്കാരൻ (kaccavaṭakkāran)}]{\ml{കച്ചവടക്കാരൻ} (kaccavaṭakkāran) — a shopkeeper, a trader}

\label{lesson:ML-C260-kaccavatakkaran}



\begin{quote}
Before the new one: say the Malayalam for chickpeas, then the Malayalam for business.
\end{quote}

\subsection*{You'll want to know: \ml{കച്ചവടക്കാരൻ}}
\textbf{\ml{കച്ചവടക്കാരൻ}} — \emph{kaccavaṭakkāran} — \textquotedblleft{}a shopkeeper, a trader\textquotedblright{}.

\begin{cousinweb}
\textbf{\ml{കച്ചവടക്കാരൻ}} (\emph{kaccavaṭakkāran}) is \emph{kaccavaṭam}, \textquotedblleft{}trade\textquotedblright{}, a word of uncertain origin, with \emph{-kkāran}, the ending for a person who does something.
\end{cousinweb}

\begin{grammarlens}[title={What you've built}]
One more word for this chapter.
\end{grammarlens}

\subsection*{Your turn}
\begin{itemize}
\raggedright
  \item \emph{Say it:} \emph{kaccavaṭakkāran}
  \item \emph{Say it:} \emph{kaccavaṭakkāran}, once more
  \item \emph{Say it:} say \emph{kaccavaṭakkāran}
  \item \emph{Recall:} say the Malayalam for chickpeas, then the Malayalam for business, then say \emph{kaccavaṭakkāran} again
\end{itemize}

\begin{tcolorbox}[breakable,colback=teal!4,colframe=teal!35!black,title={Before you move on}]
What does \ml{കച്ചവടക്കാരൻ} mean? (\textquotedblleft{}A shopkeeper, a trader\textquotedblright{}.) Say it once more.
\end{tcolorbox}
\section[\texorpdfstring{\ml{തുലാസ്} (tulāsŭ)}{തുലാസ് (tulāsŭ)}]{\ml{തുലാസ്} (tulāsŭ) — a weighing scale}

\label{lesson:ML-C260-tulasu}



\begin{quote}
Before the new one: say the Malayalam for business, then the Malayalam for a shopkeeper.
\end{quote}

\subsection*{You'll want to know: \ml{തുലാസ്}}
\textbf{\ml{തുലാസ്}} — \emph{tulāsŭ} — \textquotedblleft{}a weighing scale\textquotedblright{}.

\begin{cousinweb}
\textbf{\ml{തുലാസ്}} (\emph{tulāsŭ}) is built on Sanskrit \emph{tulā}, \textquotedblleft{}balance, weight\textquotedblright{} (CDIAL 5886); how it came by its ending is not clear.
\end{cousinweb}

\begin{grammarlens}[title={What you've built}]
One more word for this chapter.
\end{grammarlens}

\subsection*{Your turn}
\begin{itemize}
\raggedright
  \item \emph{Say it:} \emph{tulāsŭ}
  \item \emph{Say it:} \emph{tulāsŭ}, once more
  \item \emph{Say it:} say \emph{tulāsŭ}
  \item \emph{Recall:} say the Malayalam for business, then the Malayalam for a shopkeeper, then say \emph{tulāsŭ} again
\end{itemize}

\begin{tcolorbox}[breakable,colback=teal!4,colframe=teal!35!black,title={Before you move on}]
What does \ml{തുലാസ്} mean? (\textquotedblleft{}A weighing scale\textquotedblright{}.) Say it once more.
\end{tcolorbox}
\section[\texorpdfstring{\ml{കിലോ} (kilō)}{കിലോ (kilō)}]{\ml{കിലോ} (kilō) — a kilo}

\label{lesson:ML-C260-kilo}



\begin{quote}
Before the new one: say the Malayalam for a shopkeeper, then the Malayalam for a weighing scale.
\end{quote}

\subsection*{You'll want to know: \ml{കിലോ}}
\textbf{\ml{കിലോ}} — \emph{kilō} — \textquotedblleft{}a kilo\textquotedblright{}.

\begin{cousinweb}
\textbf{\ml{കിലോ}} (\emph{kilō}) $\leftarrow$ English \emph{kilo}, from French, short for \emph{kilogramme}; \emph{kilo-} is Greek \emph{chilioi}, \textquotedblleft{}thousand\textquotedblright{}.
\end{cousinweb}

\begin{grammarlens}[title={What you've built}]
One more word for this chapter.
\end{grammarlens}

\subsection*{Your turn}
\begin{itemize}
\raggedright
  \item \emph{Say it:} \emph{kilō}
  \item \emph{Say it:} \emph{kilō}, once more
  \item \emph{Say it:} say \emph{kilō}
  \item \emph{Recall:} say the Malayalam for a shopkeeper, then the Malayalam for a weighing scale, then say \emph{kilō} again
\end{itemize}

\begin{tcolorbox}[breakable,colback=teal!4,colframe=teal!35!black,title={Before you move on}]
What does \ml{കിലോ} mean? (\textquotedblleft{}A kilo\textquotedblright{}.) Say it once more.
\end{tcolorbox}
\section[\texorpdfstring{\ml{സാധനം} (sādhanaṁ)}{സാധനം (sādhanaṁ)}]{\ml{സാധനം} (sādhanaṁ) — a thing, an item, goods}

\label{lesson:ML-C260-sadhanam}



\begin{quote}
Before the new one: say the Malayalam for a weighing scale, then the Malayalam for a kilo.
\end{quote}

\subsection*{You'll want to know: \ml{സാധനം}}
\textbf{\ml{സാധനം}} — \emph{sādhanaṁ} — \textquotedblleft{}a thing, an item, goods\textquotedblright{}.

\begin{cousinweb}
\textbf{\ml{സാധനം}} (\emph{sādhanaṁ}) $\leftarrow$ Sanskrit \emph{sādhana}, \textquotedblleft{}furthering; goods\textquotedblright{}, taken over whole as a tatsama, with Malayalam's noun ending \emph{-ṁ} (CDIAL 13334).
\end{cousinweb}

\begin{grammarlens}[title={What you've built}]
One more word for this chapter.
\end{grammarlens}

\subsection*{Your turn}
\begin{itemize}
\raggedright
  \item \emph{Say it:} \emph{sādhanaṁ}
  \item \emph{Say it:} \emph{sādhanaṁ}, once more
  \item \emph{Say it:} say \emph{sādhanaṁ}
  \item \emph{Recall:} say the Malayalam for a weighing scale, then the Malayalam for a kilo, then say \emph{sādhanaṁ} again
\end{itemize}

\begin{tcolorbox}[breakable,colback=teal!4,colframe=teal!35!black,title={Before you move on}]
What does \ml{സാധനം} mean? (\textquotedblleft{}A thing, an item, goods\textquotedblright{}.) Say it once more.
\end{tcolorbox}
